\section{Related Work} \label{sec:related_work}

In this work, we leverage advances in attention-based models for meta-learning in an open-ended task space. Our agent learns a form of in-context RL algorithm, while also automatically curating the training task distribution; thus we combine two pillars of an \emph{AI generating algorithm} (AI-GA,~\citet{clune2019aiga}). The most similar work to ours is \citet{xland}, which also considers training in a vast multi-agent task space with auto-curricula and generational learning. A key difference in our work is that we focus on \emph{adaptation} (vs. zero-shot performance), and make use of large Transformer models. \citet{akkaya2019solving} also demonstrated the effectiveness of adaptive curricula while meta-learning a policy to control a robot hand, however they focused on a specific sim-to-real setting rather than a more generally capable agent. We now summarise literature related to each component of our work in turn. 

\paragraph{Procedural environment generation.} We make use of procedural content generation (PCG) to generate a vast, diverse task distribution. PCG has been studied for many years in the games community \citep{togelius2008evolving, pcg} and more recently has been used to create testbeds for RL agents \citep{pcg_illuminating, coinrun, raileanu2020ride}. Indeed, in the past few years a series of challenging PCG environments have been proposed \citep{obstacletower, cobbe2020procgen, kuettler2020nethack, samvelyan2021minihack, gym_minigrid, hafner2022benchmarking, https://doi.org/10.48550/arxiv.2206.06994}, mostly focusing on testing and improving generalisation in RL \citep{rl_generalization_survey, bhatt2022dsage, fontaine2021overcooked}. More recently there has been increased emphasis on open-ended worlds: \citet{albrecht2022avalon} proposed Avalon, a 3D world supporting complex tasks, while Minecraft \citep{johnson2016malmo} has been proposed as a challenge for Open-Endedness and RL \citep{minerl2021, fan2022minedojo, evocraft2021}, but unlike XLand it does not admit control of the full simulation stack, thereby limiting the smoothness of the task space. 

\paragraph{Open-ended learning.} A series of recent works have demonstrated the effectiveness of agent-environment co-adaptation with a distribution of tasks \citep{wang2019poet, wang2020enhanced_poet, accel}. Our approach bears resemblance to the unsupervised environment design (UED, \citet{dennis2020paired}) paradigm, since we seek to train a generalist agent without knowledge of the test tasks. One of the pioneering methods in this space was PAIRED~\citep{dennis2020paired}, which seeks to generate tasks with an RL-trained adversary. We build on Prioritised Level Replay \citep{jiang2020prioritized, jiang2021robustplr}, a method which instead curates randomly sampled environments which have high regret. Our work also relates to curriculum learning \citep{matiisen2020tscl, portelas2019teacher, sukhbaatar2018asp, openai2021asymmetric, campero2020amigo, fang2021adaptive,mu2022improving}, with the key difference that these methods typically have a specific downstream goal or task in mind. There have also been works training agents with auto-curricula over co-players, although these typically focus on singleton environments \citep{alphastar, dota} or uniformly sampled tasks \citep{Baker2020Emergent, Jaderberg859, liu2018emergent, https://doi.org/10.48550/arxiv.2203.00715}. Similar to XLand 2.0's production rule system, \citet{DBLP:conf/iclr/ZhongRG20} train agents to generalise to unobserved environment dynamics. However, they investigate zero-shot generalisation where the agent has to infer underlying environment dynamics from language descriptions, whereas AdA agents discovery these rules at test time via on-the-fly hypothesis-driven exploration over multiple trials.

\paragraph{Adaptation.} 
This work focuses on few-shot adaptation in control problems, commonly framed as \emph{meta-RL}. We focus on \emph{memory-based} meta-RL, and build upon the work of \citet{duan2017rl} and \citet{wang2016learning}, who showed that if an agent observes rewards and terminations, and the memory does not reset, a memory-based policy can implement a learning algorithm. This has proven to be an effective approach that can learn Bayes-optimal strategies~\citep{ortega2019metalearning,mikulik2020metatrained} and may have neurological analogues~\citep{wang2018prefrontal}. Indeed, our agents learn conceptual exploration strategies, something that would require the outer learner of a meta-gradient approach to estimate the return of the inner learner~\citep{stadie2018some}. Solutions in this space either rely on high-variance Monte Carlo returns~\citep{stadie2018some,garcia2019meta,vuorio2021no} or history-dependent estimators~\citep{zheng2020can}. Our work is also inspired by Alchemy~\citep{wang2102alchemy}, a meta-RL benchmark domain whose mechanics have inspired the production rules in our work. The authors use memory-based meta-RL with a small Transformer, but find that the agent's performance is only marginally better than that of a random heuristic. Transformers have also been shown to be effective for meta-RL on simple domains \citep{melo2022transformers} and for learning RL algorithms \citep{laskin2022algorithmdistillation} from offline data. Other approaches for meta-RL include meta-gradients~\citep{andrychowicz2016learning,maml_finn_icml17,xu2018meta, flennerhag2022bootstrapped}, which can be efficient but often suffer from instability and myopia~\citep{metz2021gradients,vuorio2021no,flennerhag2022bootstrapped}, and latent-variable based approaches~\citep{finn2018probabilistic,humplik2019meta,rakelly2019efficient,zintgraf2019varibad}. Adaptation also plays a critical role in robotics, with agents trained to adapt to varying terrain \citep{clavera2018learning, kumar2021rma} or damaged joints \citep{Cully2015RobotsTC}.

\paragraph{Transformers in RL and beyond.} Transformer architectures have recently shown to be highly effective for  \emph{offline} RL \citep{chen2020decision, janner2021sequence, reed2022a}, yet successes in the \emph{online} setting remain limited. One of the few works to successfully train Transformer-based policies was \citet{parisotto2020stabilizing}, who introduced several heuristics to stabilise training in a simpler, smaller-scale setting. Indeed, while we make use of a similar Transformer-XL architecture \citep{vaswani2017attention, dai2019transformer}, we demonstrate scaling laws for online meta-RL that resemble those seen in other communities, such as language \citep{DevlinCLT19, kaplan2020scaling, brown2020gpt3, rae2021scaling}. Similarly, \citet{melo2022transformers} use Transformers for fast adaptation in a smaller-scale meta-RL setting, interpreting the self-attention mechanism as a means of building an episodic memory from timestep embeddings, through the recursive application of Transformer layers. Transformer architectures have also been used in meta-learning outside of RL, for example learning general-purpose algorithms \citep{kirsch2022general} or hyperparameter optimisers \citep{optformer}. Transformers are also ubiquitous in modern large language models, which have been shown to be few-shot learners \citep{brown2020gpt3}. 
